% Generated from InterviewRecord_Group2.csv (one-off conversion; the CSV is final).
\begin{longtable}{@{}L{1.3cm} p{\dimexpr\linewidth-1.3cm-2\tabcolsep\relax}@{}}
\toprule
\textbf{ID} & \textbf{Notes} \\
\midrule
\endhead
\bottomrule
\endlastfoot
\multicolumn{2}{@{}p{\linewidth}@{}}{\textbf{G1.} What is your role in the club, and how many years have you been involved?} \\*
S2-P1 & Executive (Event Coordinator), participant, 4 years \\
S2-P2 & Coach (Structural), 3 years \\
S2-P3 & Coach (Manufacturing), 3 years \\
\addlinespace
\multicolumn{2}{@{}p{\linewidth}@{}}{\textbf{G2.} Which tools do your teams use to design, analyze, and store their bridges?} \\*
S2-P1 & Teams, SAP2000, google drive \\
S2-P2 & Autodesk inventor (model creation and safety factors analysis), SAP 200 and excel for structural analysis and member sizing optimization as well \\
S2-P3 & Design and analysis: Autodesk \\
\addlinespace
\multicolumn{2}{@{}p{\linewidth}@{}}{\textbf{G3.} What were the largest challenges you faced when you first joined the club?} \\*
S2-P1 & Primarily construction techniques and help with identifying and maintaining a timeline for the build. There was also many small decisions that are known within the club generally but when I first joined were made at whim \\
S2-P2 & The rule book and looking up past bridges made it fairly easy to design a bridge but there wasn't much guidance for manufacturing. Also lacking a dedicated space on campus with appropriate tools for the team to use later at night was hard because usually we only have time to work on the bridge in the evening. \\
S2-P3 & Designing and manufacturing joints. I think the most helpful thing we could have is a database of joinery \\
\addlinespace
\multicolumn{2}{@{}p{\linewidth}@{}}{\textbf{E1.} What information regarding the club administration is most important to consolidate?} \\*
S2-P1 & Primarily the timeline for the year, having a set timeline helps with planning events and if key milestones for teams can be consolidated, iterating potentially to improve member experience, \\
S2-P2 & Past bridges score sheets, and material sheets. \\
S2-P3 & From my perspective, where materials have been sourced, what sponsorships we have connections with. Just all around the specific places and people contacted for specific things that could be helpful for later years \\
\addlinespace
\multicolumn{2}{@{}p{\linewidth}@{}}{\textbf{E2.} When you first became an executive or coach, what were your largest difficulties?} \\*
S2-P1 & The interworking of McMaster services and working with administration for booking room and gathering supplies, also information on how to reduce costs of events \\
S2-P2 & I guess because of the new coaching structure I guess I don't really know exactly what my role is supposed to entail when it comes to how much initiative I need to take with each team or if I just wait for them to ask me questions. \\
S2-P3 & This is the first year for this coach structure, but figuring out who to talk to for what information, and how the roles have been operating in the past. The divisions feel kind of blurry, and documentation on administrative operations of any kind would help \\
\addlinespace
\multicolumn{2}{@{}p{\linewidth}@{}}{\textbf{E3.} How would you like to see the software sub team integrated into the club administration?} \\*
S2-P1 & Full integration to allow for everyone to be on top of things and understand the interworkings. \\
S2-P2 & I guess the individual team meetings as I think that would allow needs to be tailored to each team. \\
S2-P3 & I don't want them to be a separate isolated team. If that happens, they will construct stuff that might not fit the preferences and understanding of the remaining club. I do not know how this could be implemented, but if each of them has hands on design, building and competition experience it would be helpful (possibly even members that are in a team and also in the software subteam), Another option is to have a team lead that is responsible for deciding what to be added where, and that vp would be much more closely tied with the inner workings of each of the bridge building teams, meeting with coaches and other execs, etc. \\
\addlinespace
\multicolumn{2}{@{}p{\linewidth}@{}}{\textbf{C1.} What tool have you found is the most accurate and useful in predicting your bridge's performance?} \\*
S2-P1 & Not much experience doing load analysis \\
S2-P2 & Vibes we have gone with conservative guesses officially but usually once the bridge is assembled we realized our ``out there'' guess would have been better. \\
S2-P3 & Intuition and common sense: smoothness of the pieces, whether the pieces look solid or too ``dry,'' good looking proportions \\
\addlinespace
\multicolumn{2}{@{}p{\linewidth}@{}}{\textbf{C2.} How close to a model estimate do you expect your bridge's performance to be, including uncertainties?} \\*
S2-P1 & Model estimates are not ever accurate, comes down to build quality, its good for determining how good your design is theoretically \\
S2-P2 & I say new model within 25\% but eventually have the error down to 10\% \\
S2-P3 & It has never been close for me, but we have never done extensive analysis before \\
\addlinespace
\multicolumn{2}{@{}p{\linewidth}@{}}{\textbf{C3.} What are the most common reasons for bridge failures in your experience?} \\*
S2-P1 & Manufacturing issues, any weird cutting can cause major issues \\
S2-P2 & Twisting and torsion from uneven manufacturing. \\
S2-P3 & First time was concentrated stress of the force plate corner onto the bottom beam, the second was a compression member buckling \\
\addlinespace
\multicolumn{2}{@{}p{\linewidth}@{}}{\textbf{C4.} Does the club have models and recorded failure loads of past bridges that could be used to test the Toolbox?} \\*
S2-P1 & None \\
S2-P2 & You could use my teams models from the past years we still have the CAD files and the crushing videos and they were similar bridges that failed differently so that could be interesting. \\
S2-P3 & Not sure \\
\addlinespace
\multicolumn{2}{@{}p{\linewidth}@{}}{\textbf{CL1.} Is there anything about the club's information or tools that we did not ask about?} \\*
S2-P1 & No, covered most concerns \\
S2-P2 & No \\
S2-P3 & No \\
\addlinespace
\multicolumn{2}{@{}p{\linewidth}@{}}{\textbf{CL2.} Are there any items that you feel are important to include in the Toolbox that we have not discussed?} \\*
S2-P1 & A good landing page for the club, show some of the most recent best designs, good landing page for the club provides legitimacy for the club and sponsors \\
S2-P2 & I think integrating the rule book into your software could be cool to show anything to be concerned about like last year we added interior supports which made it so the car couldn't pass so if there was a dimension close to the min or max have that area flagged or something. It also could be cool for it to produce a mock score sheet. \\
S2-P3 & Above all, I want to see a nice collection of designs and models that work. Even just like a library of cad files, where when you click on it it shows its ``stats'' or something. I see this as more useful for looking at successful implementations of joints rather than entire bridges. The priority for me is the ability to quickly browse and be able to quickly see the performance scores of a bridge in competition and where it broke to quickly improve my understanding of how certain types of joints perform/how much reinforcement to plan for for a certain joint. \\
\end{longtable}
